\chapter{Les premiers jours}

La forêt engloutit rapidement les deux garçons.

Au début, Khaer Ghal ne pensa qu'à mettre le plus de distance possible entre
eux et le village. Il connaissait les bois alentour et savait dans quelle
direction avancer, mais évitait soigneusement les sentiers et les passages les
plus évidents. Si les guerriers se lançaient à leur poursuite, c'était là
qu'ils commenceraient à les chercher.

Dans l'obscurité, progresser entre les arbres ne lui posait pas de véritable
problème. Ses yeux distinguaient encore suffisamment les troncs, les rochers et
les reliefs du terrain pour qu'il puisse avancer sans lumière. Il lui fallut
un certain temps pour comprendre que Kamatsu, quelques pas derrière lui, n'avait
pas cette chance.

Il s'en aperçut lorsqu'un craquement retentit dans son dos, suivi d'un bruit
sourd. Khaer Ghal se retourna et trouva Kamatsu au sol, après qu'il eut trébuché
sur une racine.

— Tu fais beaucoup de bruit.

— Je ne vois rien.

Khaer Ghal regarda autour de lui, d'abord surpris. La forêt était sombre,
évidemment, mais il distinguait encore assez clairement ce qui l'entourait.

Puis il comprit.

Pour lui, la forêt était sombre.

Pour Kamatsu, elle était presque noire.

Allumer une torche était hors de question. Ils auraient pu être repérés de très
loin et, de toute façon, ils n'en avaient pas. Khaer Ghal revint donc vers son
ami.

— Reste près de moi.

À partir de cet instant, leur progression ralentit. Khaer Ghal avançait le
premier et prévenait Kamatsu lorsqu'une racine, une branche basse ou un trou se
trouvait devant lui. Dans les passages les plus difficiles, il lui prenait le
bras pour le guider ; ailleurs, Kamatsu posait simplement une main sur son
épaule et suivait ses mouvements.

Ils parlèrent le moins possible. Pendant plusieurs heures, il n'y eut plus que
le bruissement des feuilles, leurs respirations et le bruit prudent de leurs
pas sur le sol humide.

À mesure que la nuit avançait, pourtant, leurs pas devinrent plus lourds.
L'adrénaline de la fuite avait disparu et la fatigue commençait à prendre sa
place. Khaer Ghal continuait malgré tout à avancer, jetant régulièrement un
regard derrière eux comme s'il s'attendait encore à voir surgir les guerriers
du village entre les arbres.

— Tu crois qu'ils nous suivent ? demanda Kamatsu.

Khaer Ghal tendit l'oreille quelques secondes.

— Non.

— T'es sûr ?

Il écouta encore.

— Non.

Kamatsu soupira.

— Ça me rassure beaucoup.

Ils continuèrent encore quelque temps avant que Kamatsu ne finisse par
s'asseoir contre un arbre.

— On devrait dormir.

— Pas encore.

— Khaer Ghal, je sens plus mes jambes.

— C'est parce que t'es fatigué.

Kamatsu leva les yeux vers lui.

— Merci.

Khaer Ghal hésita. Il aurait voulu continuer, mettre encore quelques heures
entre eux et le village, mais ses propres jambes commençaient à lui donner
raison.

Kamatsu ferma les yeux.

— Moi, je dors.

Khaer Ghal ouvrit la bouche pour protester, mais un bâillement lui coupa toute
envie de défendre son argument.

\scenebreak

Trouver un endroit où dormir se révéla plus compliqué qu'il ne l'avait imaginé.
Khaer Ghal savait qu'il ne fallait pas s'installer directement au bord de
l'eau, où les animaux venaient boire et où l'humidité était plus importante.
Il savait également qu'un abri naturel valait mieux qu'un espace découvert,
qu'il fallait éviter certains creux et observer les traces alentour avant de
décider d'y passer la nuit.

Il savait beaucoup de choses. Le problème était que savoir à quoi ressemblait
un bon abri n'en faisait pas apparaître un.

Après plusieurs recherches, ils trouvèrent finalement un endroit relativement
protégé sous les racines énormes d'un vieil arbre. Elles formaient contre le
sol une cavité assez profonde pour les abriter en partie du vent. Khaer Ghal
inspecta les alentours, écarta les branches mortes et les pierres les plus
gênantes, puis ramassa des feuilles sèches et de la mousse pour les isoler un
peu du sol humide.

Il laissa Kamatsu se coucher le premier et annonça qu'il monterait la garde
quelque temps, au cas où quelque chose approcherait pendant la nuit. Assis près
de l'entrée, son épée à portée de main, il s'efforça donc de rester attentif
aux bruits inconnus de la forêt. Après les heures de marche et la fatigue
accumulée depuis leur fuite, ses bonnes intentions ne résistèrent pourtant pas
longtemps : ses paupières devinrent lourdes et il s'endormit presque aussitôt.

Lorsqu'il rouvrit les yeux, le jour filtrait déjà entre les branches. Kamatsu
dormait toujours à côté de lui et, à en juger par le calme des environs, la
forêt avait fort heureusement très bien survécu à sa première garde.

\scenebreak

Les provisions données par la mère de Khaer Ghal leur permirent de tenir les
deux premiers jours. Ils les économisèrent autant qu'ils le purent : Khaer Ghal
divisait chaque portion avec soin, rangeait ce qui restait dans le sac et
essayait d'ignorer la faim qui revenait toujours beaucoup trop vite. Kamatsu
ne protesta jamais, même lorsque les portions devinrent ridiculement petites.

Au matin du troisième jour, il ne restait plus rien.

L'eau leur posa moins de problèmes. Ils étaient restés à proximité du cours
d'eau et de ses affluents, ce qui leur permettait de remplir régulièrement leur
outre. Pour le reste, ils durent se débrouiller avec ce que la forêt acceptait
de leur donner.

Khaer Ghal n'était pas totalement démuni. Depuis son enfance, les adultes du
village lui avaient appris à reconnaître certaines plantes, quelques fruits,
les traces laissées par les animaux et plusieurs choses qu'il valait mieux ne
jamais porter à sa bouche. Il avait accompagné des chasseurs et savait observer
un terrain.

Ces connaissances suffisaient à les maintenir en vie.

Elles ne suffisaient pas à leur éviter d'avoir faim.

Le soir, lorsqu'ils s'arrêtaient, leurs ventres leur rappelaient souvent tout
ce qu'un repas chaud au village avait autrefois eu d'ordinaire.

\scenebreak

Les jours suivants se ressemblèrent sans être tout à fait identiques. Ils
marchaient pendant des heures, généralement en suivant de loin les cours d'eau
sans rester directement sur leurs berges. Khaer Ghal surveillait les traces
autour d'eux et essayait de conserver une direction ; Kamatsu apprenait à
reconnaître certains des signes que son ami lui montrait et, en échange,
remarquait parfois des choses auxquelles Khaer Ghal n'avait prêté aucune
attention.

Ils apprirent surtout à être fatigués. Leurs jambes étaient douloureuses presque
chaque soir, leurs vêtements s'accrochaient aux branches et leurs mains se
couvraient peu à peu d'éraflures. Ils découvrirent aussi qu'une nuit passée
contre des racines ne suffisait pas toujours à faire disparaître les douleurs
de la veille.

Mais ils progressaient. Chaque soir, Khaer Ghal trouvait un peu plus vite où
installer leur abri. Il distinguait mieux les bruits auxquels il devait prêter
attention de ceux qui appartenaient simplement à la forêt, et Kamatsu commençait
à savoir quelles plantes son ami acceptait de ramasser sans même avoir besoin
de demander.

Leurs repas, eux, progressaient beaucoup moins.

Un soir, Khaer Ghal rassembla ce qu'ils avaient trouvé pendant la journée et
essaya d'en faire quelque chose au-dessus d'un petit feu soigneusement
dissimulé. Le résultat n'avait rien d'exceptionnel, mais il le tendit tout de
même à Kamatsu avec une certaine fierté.

Kamatsu goûta et mâcha longuement.

— Alors ?

— C'est chaud.

— Je sais.

— C'est déjà bien.

Khaer Ghal le fixa.

— C'est mauvais ?

Kamatsu regarda ce qui restait dans son récipient.

— J'ai pas dit ça.

— Tu viens de dire que c'était chaud.

— C'est vrai.

— Kamatsu.

Un petit sourire apparut sur le visage du garçon.

— C'est très chaud.

Khaer Ghal lui lança un morceau de bois. Kamatsu l'esquiva en riant et il
essaya de conserver son air vexé quelques secondes avant de rire à son tour.
Cela faisait plusieurs jours qu'aucun des deux n'avait ri et, pendant un
instant, la fatigue, la faim et ce qu'ils avaient laissé derrière eux semblèrent
un peu moins lourds.

\scenebreak

Lorsque leur septième nuit dans la forêt arriva, les deux garçons avaient déjà
changé. Une semaine seulement s'était écoulée depuis leur départ du village,
mais elle leur semblait beaucoup plus longue. Leurs vêtements étaient sales et
déchirés par endroits, leurs visages avaient maigri, leurs mains étaient
couvertes d'éraflures et la fatigue ne disparaissait plus complètement après
une nuit de sommeil. Ils avaient presque constamment faim et avaient découvert
qu'on pouvait marcher avec les jambes douloureuses simplement parce qu'il
fallait continuer.

Pourtant, certaines choses étaient devenues plus faciles. Ils ne mettaient plus
aussi longtemps à choisir un endroit où passer la nuit. Kamatsu savait
maintenant quelles plantes Khaer Ghal reconnaissait et lesquelles il refusait
de toucher. Khaer Ghal avait cessé de dégainer son épée au moindre craquement,
et tous deux avaient appris à protéger leurs affaires de l'humidité avant de
dormir.

Ce n'était pas grand-chose. Une semaine auparavant, pourtant, ils ne savaient
pas le faire.

Ce soir-là, ils s'étaient installés sous les racines immenses d'un vieil arbre.
Un ruisseau coulait non loin et la lumière de la lune se reflétait par endroits
sur l'eau entre les troncs. Kamatsu dormait déjà lorsque Khaer Ghal resta encore
un moment assis à l'entrée de leur abri, une main près de son épée, à observer
distraitement la forêt.

Depuis sept jours, personne ne lui avait dit à quelle heure se lever, où aller,
ce qu'il devait apprendre ou quand il devait s'arrêter. Personne ne lui disait
ce qu'il devait manger, mais personne ne remplissait non plus son assiette.
Personne ne choisissait le chemin pour lui ; à lui d'apprendre à le trouver.

Cette idée lui plaisait toujours.

Khaer Ghal baissa les yeux vers le pendentif de sa mère et le prit un instant
entre ses doigts. Il avait faim. Il était fatigué. Certaines nuits, il avait
froid et, même s'il ne l'aurait probablement pas reconnu devant Kamatsu, il lui
arrivait encore d'avoir peur.

Pourtant, lorsqu'il repensait à la palissade du village, il ne souhaitait pas
y retourner.

Quelque chose remua dans les feuilles. Khaer Ghal releva la tête et écouta :
un petit animal, probablement. Quelques jours auparavant, il serait resté
éveillé pour en être certain. Cette fois, il posa son épée à portée de main,
s'allongea près de Kamatsu et ferma les yeux.

Autour d'eux, la forêt continua de vivre.

Cela faisait une semaine qu'ils étaient seuls.

Et ils étaient toujours là.